\begin{solapartat}
\figura[0.6]{sol_fig1.png}
\punts{0,75} Primer, calcularem la distància $b$: Per fer això, hem de determinar l'angle $\theta$ a partir de la llei de Snell:
\begin{align*}
n_{aire}\sin(40^\circ) = n_{aigua}\sin\theta \Rightarrow \theta &= \arcsin\left(\frac{n_{aire}\sin(40^\circ)}{n_{aigua}}\right) = \arcsin\left(\frac{1\sin(40^\circ)}{1,33}\right)\\
&= 28,9^\circ
\end{align*}
\punts{0,2} I $b$ és:
\[ \tan(\theta) = \frac{b}{2} \Rightarrow b = 2\tan(\theta) = 1,10\ \mathrm{m} \]

\punts{0,2} Fent servir la trigonometria, també obtenim $a$:
\[ \tan(40^\circ) = \frac{a}{1,6 + 0,9} \Rightarrow a = 2,5\tan(40^\circ) = 2,10\ \mathrm{m} \]
\punts{0,1} I, finalment:
\[ x = a + b = 3,20\ \mathrm{m} \]
\end{solapartat}

\begin{solapartat}
\punts{0,5} Perquè hi hagi reflexió total, l'angle d'incidència respecte de la normal ha de ser més gran o igual que l'angle límit. L'angle límit (o crític) és l'angle d'incidència respecte de la normal de la interfície pel qual l'angle de refracció és $90^\circ$. A partir de la llei de Snell:
\[ n_{aigua}\sin\theta_{lim} = n_{aire}\sin 90^\circ \]
Per tant, l'angle límit és:
\[ \theta_{lim} = \arcsin\left(\frac{n_{aire}}{n_{aigua}}\right) = \arcsin\left(\frac{1}{1,33}\right) = 48,75^\circ \]
\punts{0,25} La velocitat de propagació de la llum dins la piscina és:
\[ v = \frac{c}{n_1} = \frac{3,00\cdot 10^{8}}{1,33} = 2,26\cdot 10^{8}\ \mathrm{m/s} \]
\punts{0,25} La longitud d'ona depèn del medi, però no la freqüència; per tant, podem determinar la freqüència a partir de la velocitat de la llum i la longitud d'ona al buit:
\[ f = \frac{c}{\lambda_0} = \frac{3,00\cdot 10^{8}}{632,8\cdot 10^{-9}} = 4,74\cdot 10^{14}\ \mathrm{Hz} \]
\punts{0,25} I, finalment, la longitud d'ona dins de la piscina és:
\[ \lambda = \frac{v}{f} = \frac{2,26\cdot 10^{8}}{4,74\cdot 10^{14}} = 476\cdot 10^{-9}\ \mathrm{m} = 476\ \mathrm{nm} \]
\end{solapartat}
